\documentclass[10pt]{exam}
\usepackage[hon]{template-for-exam}
\usepackage{endnotes,graphicx}
\usepackage{tikz,graphicx,enumitem,multicol,float,my-tikz-clipart,silence,cclicenses}
\graphicspath{{./figures/}}
\WarningFilter{latex}{Label `part}
\WarningFilter{latex}{There were multiply-defined labels}

\footer{}{\thepage}{\sc \myclass}


%\def\answer#1{\footnotetext{#1}}
%\def\myquestion{\item\stepcounter{footnote}}

\def\enoteheading{}
\def\answer#1{\endnotetext{#1}}
\def\theanswers{\theendnotes}
\def\myquestion{\item\stepcounter{endnote}}

\newenvironment{myquestions}{
  \begin{multicols*}{2}\begin{enumerate}
}{
  \end{enumerate}\end{multicols*}
}
\makeatletter
\patchcmd{\enit@setresumekeys}{\def}{\gdef}{}{}
\patchcmd{\enit@setresumekeys}{\def}{\gdef}{}{}
\makeatother

\newcommand{\myheading}[1]{\uplevel{\section*{#1}}
}
\newcommand{\mysubheading}[1]{\uplevel{\subsection*{#1}}
}

%uncomment if running Pandoc
% \renewcommand{\myheading}[1]{{\section*{#1}}}
% \renewcommand{\mysubheading}[1]{\subsection*{#1}}



%%%%%%%%%%%%%%%% Dates %%%%%%%%%%%%%%%%%%%
\def\testday{Fri, Feb 13}

\def\one{Fri, Feb 27}
\def\two{Fri, Feb 27}
\def\four{Fri, Mar 6}
\def\seven{Fri, Mar 6}
%%%%%%%%%%%%%%%%%%%%%%%%%%%%%%%%%%%%%%%%%



\def\mytitle{Chapter 16A Homework (Simple Harmonic Motion)}
\title{\mytitle}
\date{\today}

\def\mymaketitle{
  \begin{flushleft}
    {\LARGE \textbf \mytitle \par}
  \end{flushleft}
}

\newcommand{\bs}[2]{\textbf{#1} (\sc #2 Points)}


\newcommand{\GFig}[2]{  
  \begin{figure}[H]
    \centering
    \includegraphics[width=#2]{11_#1_Figure.jpg}
    \caption{Giancolli Figure 11-#1.}
    \label{11-#1}
  \end{figure}
}


\begin{document}
\maketitle


\begin{tikzpicture}
  \def\len{3}
  \draw (0,0) -- (-60:\len) 
    node[circle,fill=black,text=white] (m) {$m$};
  \draw[dashed,gray] (0,0) -- (0,-\len) 
    node[circle,fill=gray,text=white] (m') {$m$};
  \draw[pattern=north east lines] 
    (-1,0) rectangle (1,0.2);
\end{tikzpicture}


\section*{Homework (collected on Mini-Test Day: \testday)}


%%%%%%%%%%%

\subsection*{Reading}

Please read the following on your own in the OpenStax textbook by the dates given.  It will give good context for class discussion.  Check off when you have completed them.

\vspace{1em}

\begin{checkboxes}
  \choice 16.1 Hooke's Law \dotfill \one
  \choice 16.2 Period \& Frequency \dotfill \two
  \choice 16.4 The Simple Pendulum \dotfill \four
  \choice 16.7 Force Oscillations \& Resonance \dotfill \seven
\end{checkboxes}


\subsection*{Problems and Conceptual Question}


Get stamps from your instructor as you complete each of the following problems.  The conceptual questions (CQ) require at least one sentence of explanation.

\vspace{1em}

\noindent
{\sc Stamps will not be given if work is not shown on a separate sheet of paper}

\vspace{1em}


\begin{tabular}{|*{2}{p{7cm}|}}
  \hline
  %
  \bs{Springs}{8}               
                    & \bs{Equations of Motion}{3} \\
  Problems \#\ref{1p-i}-\ref{1p-f}
              & Problems \#\ref{2p-i}-\ref{2p-f} \\
  Conceptual \#\ref{1q-i}-\ref{1q-f}
              &  \\[3em] \hline
  %
  \bs{Energy in SHM}{4}
                    & \bs{Simple Pendulum}{5}  \\
  Problems \#\ref{3p-i}-\ref{3p-f}
              & Problems \#\ref{4p-i}-\ref{4p-f} \\
  %Conceptual \#\ref{3q-i}-\ref{3q-f}
              & Conceptual \#\ref{4q-i}-\ref{4q-f} \\[3em] \hline
  %%
\end{tabular}

\subsection*{Bonus Problem}

\hfill

\begin{tabular}{|*{3}{p{3.5cm}|}}
  \hline
  %
  \#\ref{b1}  \\[1.5cm]
  %
  \hline
\end{tabular}



\pagebreak

\newcommand{\refsheet}{

  \section*{Equations}

  \begin{align*}
    F &= -kx &
    KE &= \frac{1}{2}mv^2 &
    PE_e &= \frac{1}{2}kx^2\\
    f_s &=\frac{1}{2\pi}\sqrt{\frac{k}{m}} &
    f_p &=\frac{1}{2\pi}\sqrt{\frac{g}{L}} &
    \omega &= 2\pi f &
    \\
    x(t) &= A \cos\left(\omega t + \phi\right) &
    v(t) &= - v_\text{max}\sin\left(\omega t + \phi\right) &
    v_\text{max} &= A\omega\\
  \end{align*}

  \section*{Selected Answers}

  \renewcommand{\enotesize}{\normalsize}
  \renewcommand{\enoteformat}{%
    \leftskip=1.8em
    \makebox[0pt][r]{\theenmark. }%
  }

  \hfill
  \theanswers
}

\refsheet


\pagebreak


\noindent
Questions and figures are from OpenStax \emph{College Physics}, 2nd ed, (Chapter 16) or Giancolli \emph{Physics: Principles with Applications}, 7th ed, (Chapter 11) unless otherwise noted.
 unless otherwise noted.



\begin{myquestions}

\myheading{Springs and Hooke's Law}
\mysubheading{Problems}


\myquestion (P4a) \label{1p-i}
	The springs of a pickup truck act like a single spring with a force constant of \SI{1.30e5}{N/m} . By how much will the truck be depressed by its maximum load of 1000~kg? \answer{7.54 cm}%; (b) \SI{3.25e4}{N/m}}
\myquestion	(P9) 
  Find the frequency of a tuning fork that takes \SI{2.50e-3}{s}  to complete one oscillation. \answer{400 Hz}
\myquestion (P18)  \label{1p-f}
  A diver on a diving board is undergoing simple harmonic motion. Her mass is 55.0 kg and the period of her motion is 0.800~s. The next diver is a male whose period of simple harmonic oscillation is 1.05~s. What is his mass if the mass of the board is negligible? \answer{94.7 kg}
% \myquestion	(Giancoli P5) 
%   A fisherman's scale stretches 3.6~cm when a 2.4-kg fish hangs from it. (a) What is the spring stiffness constant and (b) what will be the amplitude and frequency of oscillation if the fish is pulled down 2.1~cm more and released so that it oscillates up and down?

\myquestion \textbf{Bonus} \label{b1} (Giancoli P23)
  A bungee jumper with mass 65.0~kg jumps from a high bridge.  After arriving at his lowest point, he oscillates up and down, reaching a low point seven more times in 43.0~s.  He finally comes to rest 25.0~m blow the level of the bridge.  Estimate the spring stiffness constant and the unstretched length of the bungee cord assuming SHM.



\mysubheading{Conceptual Questions}

\myquestion (Giancoli P1) \label{1q-i}
  If a particle undergoes SHM with amplitude 0.21~m, what is the total distance it travels in one period? \answer{0.84 m}
% \myquestion	(CQ6) 
%   As you pass a freight truck with a trailer on a highway, you notice that its trailer is bouncing up and down slowly. Is it more likely that the trailer is heavily loaded or nearly empty? Explain your answer.
\myquestion (Giancoli CQ1) \answer{at equilibrium}
  Is the acceleration of a simple harmonic oscillator ever zero? If so, where?

\myquestion (Giancoli MC3) \answer{$4m$}
An object of mass $m$ oscillates on the end of a spring. To double the period, you should replace the object with a mass of what?

\myquestion (Giancoli MQ1) \answer{never}
  A mass on a spring in SHM (Figure~\ref{11-01}) has amplitude $A$ and period $T$. At what point in the motion (if any) is the velocity zero and the acceleration zero simultaneously?
  \GFig{01}{5cm}


\myquestion (Giancoli MC2) \answer {A, C, and D}
  An object oscillates back and forth on the end of a spring. Which of the following statements are true at some time during the course of the motion?
  \begin{parts}
    \part The object can have zero velocity and, simultaneously, nonzero acceleration.
    \part The object can have zero velocity and, simultaneously, zero acceleration.
    \part The object can have zero acceleration and, simultaneously, nonzero velocity.
    \part The object can have nonzero velocity and nonzero acceleration simultaneously.
  \end{parts}


\myquestion(Giancoli MC4)  \answer{B} \label{1q-f}
An object of mass $m$ rests on a frictionless surface and is attached to a horizontal ideal spring with spring constant $k$. The system oscillates with amplitude $A$. The oscillation frequency of this system can be increased by which of the following?
\begin{parts}
  \part decreasing $k$.
  \part decreasing $m$.
  \part increasing $A$.
  \part More than one of the above.
  \part None of the above will work.
\end{parts}

% \myquestion (CQ11) \label{1q-f}
%   How would a car bounce after a bump under each of these conditions? (a) overdamping, (b) underdamping (c) critical damping. \answer{(a) bump would hardly be absorbed; (b) car will continue to oscillate; (c) ideal case, bump will be best absorbed}

\myheading{Equations of Motion}
\mysubheading{Problems} \label{2p-i}
 
\myquestion (Giancoli P9)
  Figure~\ref{11-51} shows two examples of SHM, labeled A and B. For each, what is (a) the amplitude, (b) the frequency, and (c) the period?
  \GFig{51}{5.5cm} \answer{2.25 m \& 3.5 m; 0.25 Hz \& 0.5 Hz; 4 s \& 2 s}

\myquestion (Giancoli P13) \label{2p-f}
  A 1.65-kg mass stretches a vertical spring 0.215~m. If the spring is stretched an additional 0.130~m and released, how long does it take to reach the (new) equilibrium position again? \answer{0.23 sec}

\myquestion (Giancoli P8) 
  A vertical spring with spring stiffness constant 305~N/m oscillates with an amplitude of 28.0~cm when 0.235~kg hangs from it. The mass passes through the equilibrium point ($y=0$) with positive velocity at $t=0$, which means that $\phi=3\pi/2$. (a) What equation describes this motion as a function of time? (b) At what times will the spring be longest and shortest? \answer{$x = 0.28\cos (36.0t + 3\pi/2)$; $t = 0.04$ s \& 0.13 s}

\myheading{Energy in SHM}
\mysubheading{Problems}


\myquestion (Giancoli P14) \label{3p-i}
  A 1.15-kg mass oscillates according to the equation
  %
  \begin{align*}
    x&=0.650\cos(8.4t)
  \end{align*}
  %
  where $x$ is in meters and $t$ in seconds. Determine (a) the amplitude, (b) the frequency, (c) the total energy, and (d) the kinetic energy and potential energy when $x=0.36$~m. \answer{0.650 m;  1.34 Hz;  Total $ = 17.2$ J; $PE = 5.26$ J; $KE = 11.92$ J}

\myquestion (Giancoli P20)  \answer{5.5 cm;  0.59 m/s}
  An object with mass 2.7~kg is executing simple harmonic motion, attached to a spring with spring constant $k=310$~N/m. When the object is 0.020~m from its equilibrium position, it is moving with a speed of 0.55~m/s. (a) Calculate the amplitude of the motion. (b) Calculate the maximum speed attained by the object.

\myquestion (P38) \answer{(a) 0.335 m/s; (b) \SI{5.61e-4}{J}} \label{3p-f}
  A novelty clock has a 0.0100-kg mass object bouncing on a spring that has a force constant of 1.25~N/m. What is the maximum velocity of the object if the object bounces 3.00~cm above and below its equilibrium position? (b) How many joules of kinetic energy does the object have at its maximum velocity?

% \mysubheading{Conceptual Questions}
% \myquestion (Giancoli P11) \label{3q-i}
%   At what displacement of a Simple Harmonic Oscillator is the energy half kinetic and half potential?
% \myquestion (Giancoli P17) \label{3q-f}
%   If one oscillation has 3.0 times the energy of a second one of equal frequency and mass, what is the ratio of their amplitudes?


\myheading{Simple Pendulum}
\mysubheading{Problems}


\myquestion (Giancoli P27)  \answer{3.0 sec} \label{4p-i}
  A pendulum has a period of 1.85~s on Earth. What is its period on Mars, where the acceleration of gravity is about \SI{3.63}{m/s^2}?

\myquestion (Giancoli P28) \answer{0.994 m}
  How long must a simple pendulum be if it is to make exactly one swing per second? (That is, one complete oscillation takes exactly 2.0~s.)

\myquestion (Giancoli P29) \answer{1.8 sec;  0.56 Hz} \label{4p-f}
  A pendulum makes 28 oscillations in exactly 50~s. What is its (a) period and (b) frequency?

% \myquestion (P25) 
%   How long does it take a child on a swing to complete one swing if her center of gravity is 4.00~m below the pivot?



\mysubheading{Conceptual Questions}

\myquestion (CQ8) \answer{lengthen} \label{4q-i}
  Pendulum clocks are made to run at the correct rate by adjusting the pendulum's length. Suppose you move from one city to another where the acceleration due to gravity is slightly greater, taking your pendulum clock with you, will you have to lengthen or shorten the pendulum to keep the correct time, other factors remaining constant? Explain your answer.

\myquestion (P30) \answer{period gets cut by $1/\sqrt{2}$} \label{4q-f}
  What is the effect on the period of a pendulum if you double its length?


\end{myquestions}

\end{document}